%%
%%
\section{File \texttt{input.zone}}
\subsection{Purpose}
This file is the main file controlling the \swZone executable. Its purpose is to create spatially distributed (cell- or edge-based) parameter files to be read as inputs by \swSolver. The parameters are assumed to be uniform within zones defined by polygons.
%
\subsection{Format description}
Empty lines or lines starting with hash '\#' or equal '=' signs are ignored. The spelling of the parameters is case sensitive.

This text file contains four mandatory blocks that are delimited by specific line strings.

\begin{scriptsize}
\begin{tabular}{l | l | l}
\hline
Block & starting line string & ending line string \\
\hline
Block 1 & \texttt{Block1: zones list BEGIN (compulsory line)} & \texttt{Block1: zones list END (compulsory line)} \\
Block 2 & \texttt{Block2: default parameter value BEGIN (compulsory line)} & \texttt{Block2: default parameter value END (compulsory line)} \\
Block 3 & \texttt{Block3: linking variable-zone BEGIN (compulsory line)} & \texttt{Block3: linking variable-zone END (compulsory line)} \\
Block 4 & \texttt{Block4: output file names BEGIN (compulsory line)} & \texttt{Block4: output file name END (compulsory line)} \\
\hline
\end{tabular}
\end{scriptsize}

\subsubsection{Block 1}
Each line of Block 1 defines one polygon. For each polygon, the summits can be either identified by their coordinates or as a node of the mesh. Each line should correspond to one of the two possible formats:
\begin{itemize}
    \item 
    \begin{lstlisting}
    <NAME> $P$ $x_1$ $y_1$ $x_2$ $y_2$ ... $x_i$ $y_i$ ... \end{lstlisting}
where \texttt{<NAME>} is the name of the zone, $P$ indicates that the summits are defined by physical coordinates, and the $x_i,y_i$ couples are the coordinates of the polynomial vertices.
    \item 
        \begin{lstlisting}
<NAME> $N$ $n_1$ $n_2$ ... $n_i$ ...
\end{lstlisting}
where \texttt{<NAME>} is the name of the zone, $N$ indicates that the summits are defined by mesh nodes, and the $n_i$ are the node identifiers for the vertices of the polynomial. The $n_i$ must therefore correspond to actual mesh node numbers.
\end{itemize}


\subsubsection{Block 2}
Each line of Block 2 defines the default value for each parameter. The format is the following
\begin{lstlisting}
<PARAM> <DEFAULT VALUE>
\end{lstlisting}
where \texttt{<PARAM>} is the name of the parameter and \texttt{<DEFAULT VALUE>} is the default value for this parameter. Note that
\begin{itemize}
\item SI units are used for all parameters
\item \swZone does not check whether the value is physically meaningful
\end{itemize}. 

\subsubsection{Block 3}
Each line in Block 3 assigns a numerical value to a parameter within a polygon previously defined in Block 1. The format is the following
\begin{lstlisting}
<PARAM> <VALUE> <NAME>
\end{lstlisting}
where \texttt{<PARAM>} is the name of the parameter, \texttt{<VALUE>} is its numerical value and \texttt{<NAME>} is the name of one of the polygons defined in Block 1. 

The parameters assigned in this Block 3 must have a default value defined in Block 2.

If two or more blocks overlap, the parameter values provided in the last polygon of the list will overwrite those previously defined in the overlapping area. 

\subsubsection{Block 4}
Each line in Block 4 defines the files to be created by \swZone and the parameters to be written into them. The format is the following
\begin{lstlisting}
<FILENAME> <PARAM_1> ... <PARAM_i> ...
\end{lstlisting}
where \texttt{<FILENAME>} is the output file name and \texttt{<PARAM\_i>} are the parameters to be written into this file. 

\subsection{Sample file}
In this sample file, the parameter file \texttt{initial\_conditions\_map.txt} is to be filled with three different parameters : h, q and v.
The default values for h, q and v are 50, 0.0 and -1.0 respectively.
A quadrangular polygon (called "Zone1") is defined by the coordinates of its vertices $(-1,-1), (400,-1),(400,1), (-1,1)$. Within this polygon, the parameter \texttt{h} is set to 10, while q and v are equal to their default values. 
A second polygon (called "Zone2") is defined as a draft (the \# character indicates that this is a comment line) using the Ids of its 7 vertices. 

\begin{verbatim}
= line of comments
# symbol to comment : (# =)
Block1: zones list BEGIN (compulsory line)
Zone1	P	-1	-1	400	-1	400	1	-1	1
#Zone2	N	1   2   3   4   5   8   12
Block1: zones list END (compulsory line)
Block2: default parameter value BEGIN (compulsory line)
h	50.
q	0.
v	-1.
Block2: default parameter value END (compulsory line)
Block3: linking variable-zone BEGIN (compulsory line)
h	10.	Zone1
Block3: linking variable-zone END (compulsory line)
Block4: output file name BEGIN (compulsory line)
initial_conditions_map.txt  h   q   v
Block4: output file name END (compulsory line)
\end{verbatim}